\documentclass{birkau}

\usepackage{amsmath,amssymb,url}
\usepackage{hyperref}
\usepackage[all,cmtip]{xy}
\usepackage{enumitem}
\usepackage{mathtools}
\input diagxy

\theoremstyle{plain}
 \newtheorem{thm}{Theorem}[section]
 \newtheorem{lem}[thm]{Lemma}
 \newtheorem{prop}[thm]{Proposition}
 \newtheorem{cor}[thm]{Corollary}
\theoremstyle{definition}
 \newtheorem{defn}[thm]{Definition}
 \newtheorem{exmp}[thm]{Example}
 \newtheorem{rem}[thm]{Remark}
 \newtheorem{note}[thm]{Note}
 \newtheorem{ques}[thm]{Question}
 \newtheorem{case}[thm]{Case}


\DeclareMathAlphabet{\mathcal}{OMS}{cmsy}{m}{n}
\DeclareMathOperator*{\colim}{colim}
\DeclareMathOperator{\Cone}{Cone}
\DeclareMathOperator{\dom}{dom}
\DeclareMathOperator{\cod}{cod}
\DeclareMathOperator{\Lan}{Lan}
\DeclareMathOperator{\Ran}{Ran}
\DeclareMathOperator{\ev}{ev}
\DeclareMathOperator{\ub}{ub}
\DeclareMathOperator{\lb}{lb}
\DeclareMathOperator{\id}{id}
\DeclareMathOperator{\ob}{ob}
\DeclareMathOperator{\supp}{supp}
\DeclareMathOperator{\sIm}{\mathsf{Im}}
\DeclareMathOperator{\Int}{Int}

\def\oto{{\bfig\morphism<180,0>[\mkern-4mu`\mkern-4mu;]\place(86,0)[\circ]\efig}}
\def\rto{{\bfig\morphism<180,0>[\mkern-4mu`\mkern-4mu;]\place(82,0)[\mapstochar]\efig}}
\def\nra{\relbar\joinrel\joinrel\mapstochar\joinrel\joinrel\rightarrow}
\def\nla{\leftarrow\joinrel\joinrel\joinrel\mapstochar\joinrel\relbar}
%\def\ldd{\mkern1mu{\bfig\morphism/{*}-/<90,90>[`;]\efig}\mkern-7mu}
%\def\rdd{\mkern-7mu{\bfig\morphism/-{*}/<90,-90>[`;]\efig}\mkern1mu}

\newcommand{\da}{\downarrow}
\newcommand{\ua}{\uparrow}
\newcommand{\ra}{\rightarrow}
\newcommand{\la}{\leftarrow}
\newcommand{\Ra}{\Rightarrow}
\newcommand{\La}{\Leftarrow}
\newcommand{\lra}{\longrightarrow}
\newcommand{\lda}{\swarrow}
\newcommand{\rda}{\searrow}
\newcommand{\bs}{\backslash}
\newcommand{\Lra}{\Longrightarrow}
\newcommand{\Lla}{\Longleftarrow}
\newcommand{\rat}{\rightarrowtail}
\newcommand{\bv}{\bigvee}
\newcommand{\bw}{\bigwedge}
\newcommand{\tl}{\triangleleft}
\newcommand{\tr}{\triangleright}
\newcommand{\vd}{\vdash}
\newcommand{\dv}{\dashv}
\newcommand{\rhu}{\rightharpoonup}
\newcommand{\nat}{\natural}
\newcommand{\ola}{\overleftarrow}
\newcommand{\ora}{\overrightarrow}
\newcommand{\rqa}{\rightsquigarrow}
\newcommand{\lar}{\looparrowright}
\newcommand{\od}{\odot}
\newcommand{\opl}{\oplus}
\newcommand{\dbv}{\displaystyle\bv}
\newcommand{\dbw}{\displaystyle\bw}

\renewcommand{\phi}{\varphi}
\newcommand{\al}{\alpha}
\newcommand{\be}{\beta}
\newcommand{\up}{\upsilon}
\newcommand{\Up}{\Upsilon}
\newcommand{\ep}{\varepsilon}
\newcommand{\de}{\delta}
\newcommand{\De}{\Delta}
\newcommand{\ga}{\gamma}
\newcommand{\Ga}{\Gamma}
\newcommand{\Lam}{\Lambda}
\newcommand{\lam}{\lambda}
\newcommand{\si}{\sigma}
\newcommand{\Si}{\Sigma}
\newcommand{\Om}{\Omega}
\newcommand{\om}{\omega}
\newcommand{\ka}{\kappa}

\newcommand{\CA}{\mathcal{A}}
\newcommand{\CB}{\mathcal{B}}
\newcommand{\CC}{\mathcal{C}}
\newcommand{\CD}{\mathcal{D}}
\newcommand{\CE}{\mathcal{E}}
\newcommand{\CF}{\mathcal{F}}
\newcommand{\CG}{\mathcal{G}}
\newcommand{\CH}{\mathcal{H}}
\newcommand{\CI}{\mathcal{I}}
\newcommand{\CJ}{\mathcal{J}}
\newcommand{\CK}{\mathcal{K}}
\newcommand{\CL}{\mathcal{L}}
\newcommand{\CM}{\mathcal{M}}
\newcommand{\CN}{\mathcal{N}}
\newcommand{\CO}{\mathcal{O}}
\newcommand{\CP}{\mathcal{P}}
\newcommand{\CQ}{\mathcal{Q}}
\newcommand{\CR}{\mathcal{R}}
\newcommand{\CS}{\mathcal{S}}
\newcommand{\CT}{\mathcal{T}}
\newcommand{\CU}{\mathcal{U}}
\newcommand{\CV}{\mathcal{V}}
\newcommand{\CW}{\mathcal{W}}

\newcommand{\sA}{\mathsf{A}}
\newcommand{\sB}{\mathsf{B}}
\newcommand{\sC}{\mathsf{C}}
\newcommand{\sD}{\mathsf{D}}
\newcommand{\sE}{\mathsf{E}}
\newcommand{\sF}{\mathsf{F}}
\newcommand{\sG}{\mathsf{G}}
\newcommand{\sH}{\mathsf{H}}
\newcommand{\sI}{\mathsf{I}}
\newcommand{\sJ}{\mathsf{J}}
\newcommand{\sK}{\mathsf{K}}
\newcommand{\sL}{\mathsf{L}}
\newcommand{\sM}{\mathsf{M}}
\newcommand{\sN}{\mathsf{N}}
\newcommand{\sO}{\mathsf{O}}
\newcommand{\sP}{\mathsf{P}}
\newcommand{\sQ}{\mathsf{Q}}
\newcommand{\sR}{\mathsf{R}}
\newcommand{\sS}{\mathsf{S}}
\newcommand{\sT}{\mathsf{T}}
\newcommand{\sU}{\mathsf{U}}
\newcommand{\sV}{\mathsf{V}}
\newcommand{\sW}{\mathsf{W}}
\newcommand{\sX}{\mathsf{X}}
\newcommand{\sY}{\mathsf{Y}}
\newcommand{\sZ}{\mathsf{Z}}
\newcommand{\se}{\mathsf{e}}
\newcommand{\sm}{\mathsf{m}}
\newcommand{\sd}{\mathsf{d}}
\newcommand{\sk}{\mathsf{k}}
\newcommand{\so}{\mathsf{o}}
\newcommand{\sq}{\mathsf{q}}
\newcommand{\sy}{\mathsf{y}}
\newcommand{\sfs}{\mathsf{s}}
\newcommand{\Fix}{\mathsf{Fix}}
\newcommand{\cl}{\mathsf{cl}}

\newcommand{\bbA}{\mathbb{A}}
\newcommand{\bbB}{\mathbb{B}}
\newcommand{\bbC}{\mathbb{C}}
\newcommand{\bbD}{\mathbb{D}}
\newcommand{\bbE}{\mathbb{E}}
\newcommand{\bbH}{\mathbb{H}}
\newcommand{\bbI}{\mathbb{I}}
\newcommand{\bbL}{\mathbb{L}}
\newcommand{\bbN}{\mathbb{N}}
\newcommand{\bbP}{\mathbb{P}}
\newcommand{\bbQ}{\mathbb{Q}}
\newcommand{\bbR}{\mathbb{R}}
\newcommand{\bbS}{\mathbb{S}}
\newcommand{\bbT}{\mathbb{T}}
\newcommand{\bbU}{\mathbb{U}}
\newcommand{\bbW}{\mathbb{W}}
\newcommand{\bbX}{\mathbb{X}}
\newcommand{\bbY}{\mathbb{Y}}
\newcommand{\bbZ}{\mathbb{Z}}

\newcommand{\BB}{\mathbf{B}}
\newcommand{\BD}{\mathbf{D}}
\newcommand{\BH}{\mathbf{H}}
\newcommand{\BK}{\mathbf{K}}
\newcommand{\Bf}{\mathbf{f}}
\newcommand{\Bg}{\mathbf{g}}
\newcommand{\Bh}{\mathbf{h}}
\newcommand{\Bj}{\mathbf{j}}
\newcommand{\Bu}{\mathbf{u}}
\newcommand{\Bv}{\mathbf{v}}

\newcommand{\FA}{\mathfrak{A}}
\newcommand{\FB}{\mathfrak{B}}
\newcommand{\FC}{\mathfrak{C}}
\newcommand{\FD}{\mathfrak{D}}
\newcommand{\FE}{\mathfrak{E}}
\newcommand{\FF}{\mathfrak{F}}
\newcommand{\FG}{\mathfrak{G}}
\newcommand{\FH}{\mathfrak{H}}
\newcommand{\FI}{\mathfrak{I}}
\newcommand{\FK}{\mathfrak{K}}
\newcommand{\FL}{\mathfrak{L}}
\newcommand{\FM}{\mathfrak{M}}
\newcommand{\FP}{\mathfrak{P}}
\newcommand{\FQ}{\mathfrak{Q}}
\newcommand{\FR}{\mathfrak{R}}
\newcommand{\FS}{\mathfrak{S}}
\newcommand{\FT}{\mathfrak{T}}
\newcommand{\FU}{\mathfrak{U}}
\newcommand{\FV}{\mathfrak{V}}
\newcommand{\Fb}{\mathfrak{b}}
\newcommand{\Fi}{\mathfrak{i}}
\newcommand{\Fj}{\mathfrak{j}}
\newcommand{\Fr}{\mathfrak{r}}
\newcommand{\Fs}{\mathfrak{s}}
\newcommand{\Ff}{\mathfrak{f}}
\newcommand{\Fx}{\mathfrak{x}}
\newcommand{\Fy}{\mathfrak{y}}
\newcommand{\Fz}{\mathfrak{z}}

\newcommand{\QB}{\CQ_{\CB}}
\newcommand{\Alg}{\mathbf{Alg}}
\newcommand{\Arr}{\mathbf{Arr}}
\newcommand{\Cat}{\mathbf{Cat}}
\newcommand{\Cls}{\mathbf{Cls}}
\newcommand{\ClsDist}{\mathbf{ClsDist}}
\newcommand{\ClsCloRel}{\mathbf{ClsRel}_{\cl}}
\newcommand{\CAT}{\mathbf{CAT}}
\newcommand{\CCat}{\mathbf{CCat}}
\newcommand{\Chu}{\mathbf{Chu}}
\newcommand{\ChuCon}{\mathbf{ChuCon}}
\newcommand{\Dis}{\mathbf{Dis}}
\newcommand{\Dist}{\mathbf{Dist}}
\newcommand{\Inf}{\mathbf{Inf}}
\newcommand{\Info}{\mathbf{Info}}
\newcommand{\IntDist}{\mathbf{IntDist}}
\newcommand{\IntRel}{\mathbf{IntRel}}
\newcommand{\TopRel}{\mathbf{TopRel}}
\newcommand{\Map}{\mathbf{Map}}
\newcommand{\Mat}{\mathbf{Mat}}
\newcommand{\Met}{\mathbf{Met}}
\newcommand{\Mod}{\mathbf{Mod}}
\newcommand{\Mon}{\mathbf{Mon}}
\newcommand{\GMon}{\mathbf{GMon}}
\newcommand{\Ord}{\mathbf{Ord}}
\newcommand{\Quant}{\mathbf{Quant}}
\newcommand{\Rel}{\mathbf{Rel}}
\newcommand{\Set}{\mathbf{Set}}
\newcommand{\CcSet}{\mathbf{CcSet}}
\newcommand{\GSet}{\mathbf{GSet}}
\newcommand{\Sup}{\mathbf{Sup}}
\newcommand{\ParApp}{\mathbf{ParApp}}
\newcommand{\ParMet}{\mathbf{ParMet}}
\newcommand{\ParOrd}{\mathbf{ParOrd}}
\newcommand{\ParTop}{\mathbf{ParTop}}
\newcommand{\ProbMet}{\mathbf{ProbMet}}
\newcommand{\ProbParMet}{\mathbf{ProbParMet}}
\newcommand{\ParMetMon}{\mathbf{ParMetMon}}
\newcommand{\UMet}{\mathbf{UMet}}
\newcommand{\QApp}{\CQ\text{-}\App}
\newcommand{\QsApp}{\CQ\text{-}\sApp}
\newcommand{\QCat}{\CQ\text{-}\Cat}
\newcommand{\QSymCat}{\CQ\text{-}\mathbf{SymCat}}
\newcommand{\QCcSymCat}{\CQ\text{-}\mathbf{CcSymCat}_{\sf sk}}
\newcommand{\QMon}{\sQ\text{-}\Mon}
\newcommand{\DQMon}{\DQ\text{-}\Mon}
\newcommand{\QGMon}{\sQ\text{-}\GMon}
\newcommand{\QSet}{\sQ\text{-}\Set}
\newcommand{\QGSet}{\sQ\text{-}\GSet}
\newcommand{\RCat}{\CR\text{-}\Cat}
\newcommand{\sQCat}{\sQ\text{-}\Cat}
\newcommand{\QCAT}{\CQ\text{-}\CAT}
\newcommand{\QCCat}{\CQ\text{-}\CCat}
\newcommand{\QCls}{\CQ\text{-}\Cls}
\newcommand{\QDist}{\CQ\text{-}\Dist}
\newcommand{\QMap}{\sQ\text{-}\mathbf{Map}}
\newcommand{\QParMap}{\sQ\text{-}\Map^{\partial}}
\newcommand{\QSymDist}{\CQ\text{-}\mathbf{SymDist}}
\newcommand{\QCcSymDist}{\CQ\text{-}\mathbf{CcSymDist}}
\newcommand{\Rd}{\mathbf{Rd}}
\newcommand{\Rc}{\mathbf{Rc}}
\newcommand{\QChu}{\CQ\text{-}\Chu}
\newcommand{\QInf}{\CQ\text{-}\Inf}
\newcommand{\QInfo}{\CQ\text{-}\Info}
\newcommand{\QInt}{\CQ\text{-}\Int}
\newcommand{\QIntDist}{\CQ\text{-}\IntDist}
\newcommand{\QIntRel}{\CQ\text{-}\IntRel}
\newcommand{\QClsDist}{\CQ\text{-}\ClsDist}
\newcommand{\QClsCloDist}{(\QClsDist)_{\cl}}
\newcommand{\sQMap}{\{\star\}/\QMap}
\newcommand{\QMat}{\CQ\text{-}\Mat}
\newcommand{\QMod}{\CQ\text{-}\Mod}
\newcommand{\QOrd}{\sQ\text{-}\Ord}
\newcommand{\QRel}{\sQ\text{-}\Rel}
\newcommand{\QSup}{\CQ\text{-}\Sup}
\newcommand{\VCat}{\CV\text{-}\Cat}
\newcommand{\VCAT}{\CV\text{-}\CAT}
\newcommand{\VDis}{\CV\text{-}\Dis}
\newcommand{\VDist}{\CV\text{-}\Dist}
\newcommand{\VMat}{\CV\text{-}\Mat}
\newcommand{\VMod}{\CV\text{-}\Mod}
\newcommand{\VRel}{\CV\text{-}\Rel}
\newcommand{\TopQ}{\Top(\CQ)}
\newcommand{\ClsQ}{\Cls(\CQ)}
\newcommand{\ClsbQ}{\Cls_{\bot}(\CQ)}
\newcommand{\RSet}{[0,1]_*\text{-}\Set}
\newcommand{\RCcSet}{[0,1]_*\text{-}\CcSet}
\newcommand{\RSetM}{[0,1]_*\text{-}\underline{\Set}}
\newcommand{\RCcSetM}{[0,1]_*\text{-}\underline{\CcSet}}
\newcommand{\RDist}{[0,1]_*\text{-}\mathbf{Dist}}

\newcommand{\aR}{R^{\forall}}
\newcommand{\eR}{R_{\exists}}
\newcommand{\dR}{R^{\da}}
\newcommand{\uR}{R_{\ua}}
\newcommand{\aS}{S^{\forall}}
\newcommand{\eS}{S_{\exists}}
\newcommand{\dS}{S^{\da}}
\newcommand{\uS}{S_{\ua}}
\newcommand{\hga}{\widehat{\ga}}
\newcommand{\tga}{\widetilde{\ga}}
\newcommand{\aphi}{\phi^{\forall}}
\newcommand{\ephi}{\phi_{\exists}}
\newcommand{\dphi}{\phi^{\da}}
\newcommand{\uphi}{\phi_{\ua}}
\newcommand{\apsi}{\psi^{\forall}}
\newcommand{\epsi}{\psi_{\exists}}
\newcommand{\dpsi}{\psi^{\da}}
\newcommand{\upsi}{\psi_{\ua}}
\newcommand{\hphi}{\widehat{\phi}}
\newcommand{\tphi}{\widetilde{\phi}}
\newcommand{\hpsi}{\widehat{\psi}}
\newcommand{\tpsi}{\widetilde{\psi}}
\newcommand{\hzeta}{\widehat{\zeta}}
\newcommand{\tzeta}{\widetilde{\zeta}}
\newcommand{\heta}{\widehat{\eta}}
\newcommand{\teta}{\widetilde{\eta}}
\newcommand{\hxi}{\widehat{\xi}}
\newcommand{\txi}{\widetilde{\xi}}
\newcommand{\htheta}{\widehat{\theta}}
\newcommand{\ttheta}{\widetilde{\theta}}
\newcommand{\hf}{\hat{f}}
\newcommand{\hF}{\hat{F}}
\newcommand{\hK}{\hat{\sK}}
\newcommand{\hI}{\hat{\sI}}
\newcommand{\hKd}{\hK_{\sd}}
\newcommand{\hId}{\hI_{\sd}}
\newcommand{\hIo}{\hI_{\so}}

\newcommand{\lphi}{\ola{\phi}}
\newcommand{\lpsi}{\ola{\psi}}
\newcommand{\lzeta}{\ola{\zeta}}
\newcommand{\rphi}{\ora{\phi}}
\newcommand{\rpsi}{\ora{\psi}}
\newcommand{\rzeta}{\ora{\zeta}}
\newcommand{\raf}{\ora{f}}
\newcommand{\rag}{\ora{g}}
\newcommand{\laf}{\ola{f}}
\newcommand{\lag}{\ola{g}}
\newcommand{\tf}{\widetilde{f}}

\newcommand{\sCd}{\sC^{\dag}}
\newcommand{\sPd}{\sP^{\dag}}
\newcommand{\syd}{\sy^{\dag}}
\newcommand{\sYd}{\sY^{\dag}}
\newcommand{\Yd}{Y^{\dag}}
\newcommand{\co}{{\rm co}}
\newcommand{\op}{{\rm op}}
\newcommand{\PA}{\sP\bbA}
\newcommand{\PB}{\sP\bbB}
\newcommand{\PC}{\sP\bbC}
\newcommand{\PX}{\sP X}
\newcommand{\PY}{\sP Y}
\newcommand{\PZ}{\sP Z}
\newcommand{\CX}{\sC X}
\newcommand{\CY}{\sC Y}
\newcommand{\CZ}{\sC Z}
\newcommand{\SX}{\sS X}
\newcommand{\SY}{\sS Y}
\newcommand{\FX}{\mathfrak{X}}
\newcommand{\FY}{\mathfrak{Y}}
\newcommand{\FYd}{\FY^{\dag}}
\newcommand{\Fyd}{\Fy^{\dag}}
\newcommand{\FZ}{\mathfrak{Z}}
\newcommand{\LX}{\sL X}
\newcommand{\LY}{\sL Y}
\newcommand{\LZ}{\sL Z}
\newcommand{\PdA}{\sPd\bbA}
\newcommand{\PdB}{\sPd\bbB}
\newcommand{\PdC}{\sPd\bbC}
\newcommand{\PdX}{\sPd X}
\newcommand{\PdY}{\sPd Y}
\newcommand{\PdZ}{\sPd Z}
\newcommand{\CdX}{\sCd X}
\newcommand{\CdY}{\sCd Y}
\newcommand{\CdZ}{\sCd Z}
\newcommand{\Mphi}{\sM\phi}
\newcommand{\Mpsi}{\sM\psi}
\newcommand{\Kphi}{\sK\phi}
\newcommand{\Kpsi}{\sK\psi}

\newcommand{\ds}{\displaystyle}
\newcommand{\BQ}{\BB(\CQ)}
\newcommand{\BsQ}{\BB(\sQ)}
\newcommand{\KsQ}{\BB_*(\sQ)}
\newcommand{\KssQ}{\BB_*^{\circ}(\sQ)}
\newcommand{\KC}{\BB_*(C_3)}
\newcommand{\BL}{\BB(L)}
\newcommand{\BsQCat}{\BsQ\text{-}\Cat}
\newcommand{\DQ}{\sD\sQ}
\newcommand{\DsQ}{\BD(\sQ)}
\newcommand{\HsQ}{\BD_*(\sQ)}
\newcommand{\HssQ}{\BD_*(\sQ)^{\circ}}
\newcommand{\HC}{\BD_*(C_3)}
\newcommand{\DL}{\BD(L)}
\newcommand{\DsQCat}{\DsQ\text{-}\Cat}
\newcommand{\ArrQ}{\Arr(\CQ)}
\newcommand{\ChuConQ}{\ChuCon(\CQ)}
\newcommand{\DDe}{\sD\De}
\newcommand{\FDist}{\mathbf{FDist}}
\newcommand{\FOrd}{\mathbf{FOrd}}
\newcommand{\FRel}{\mathbf{FRel}}
\newcommand{\FSup}{\mathbf{FSup}}
\newcommand{\POrd}{\mathbf{POrd}}
\newcommand{\QFDist}{\FQ\text{-}\FDist}
\newcommand{\QFOrd}{\FQ\text{-}\FOrd}
\newcommand{\QFRel}{\FQ\text{-}\FRel}
\newcommand{\QFSup}{\FQ\text{-}\FSup}
\newcommand{\QPOrd}{\FQ\text{-}\POrd}
\newcommand{\DQCat}{\DQ\text{-}\Cat}
\newcommand{\DQMat}{\DQ\text{-}\Mat}
\newcommand{\DQRel}{\DQ\text{-}\Rel}
\newcommand{\pid}{\pi^{\dag}}

\newcommand{\PDX}{\mathsf{PD}(X)}
\newcommand{\PCX}{\mathsf{PC}(X)}
\newcommand{\ldd}{\mathrel{/}}
\newcommand{\rdd}{\mathrel{\bs}}
\newcommand{\with}{\mathrel{\&}}
\newcommand{\Gb}{\FG_{\Fb}}
\newcommand{\Gf}{\FG_{\Ff}}
\newcommand{\CGb}{\CG_{\Fb}}
\newcommand{\CGf}{\CG_{\Ff}}
\newcommand{\sGb}{\sG_{\Fb}}
\newcommand{\sGf}{\sG_{\Ff}}
\newcommand{\Lb}{\FL_{\Fb}}
\newcommand{\Lf}{\FL_{\Ff}}
\newcommand{\CLb}{\CL_{\Fb}}
\newcommand{\CLf}{\CL_{\Ff}}
\newcommand{\sLb}{\sL_{\Fb}}
\newcommand{\sLf}{\sL_{\Ff}}
\renewcommand{\leq}{\leqslant}
\renewcommand{\geq}{\geqslant}
\newcommand{\sQCats}{\sQCat_{\sfs}}
\newcommand{\DsQCats}{\DsQCat_{\sfs}}
\newcommand{\BsQCats}{\BsQCat_{\sfs}}
\newcommand{\RQDist}{\BD(\QDist)_{\rm reg}}
\newcommand{\QCD}{(\QSup)_{\rm ccd}}
\newcommand{\ts}{\mathrel{\otimes}}

\numberwithin{equation}{section}
\allowdisplaybreaks

\renewcommand{\qedsymbol}{\ensuremath{\blacksquare}}

\begin{document}

\title{Cartesian closedness of the category of real-valued sets, II}

\author[L. Shen]{Lili Shen}
\address{School of Mathematics, Sichuan University, Chengdu 610064, China}
\email{shenlili@scu.edu.cn}

\corrauthor[J. Zhang]{Jian Zhang}
\address{School of Mathematics, Sichuan University, Chengdu 610064, China}
\email{zhangjian.2025@qq.com}

\subjclass{03E72, 18D15, 18F75}

\keywords{Quantales, Quantaloids, Real-valued sets, Left-continuous t-norms, Cartesian closed categories}

\begin{abstract}
Let $[0,1]_*$ be the unit interval $[0,1]$ equipped with a left-continuous t-norm $*$. We prove that $[0,1]_*\text{-}\mathbf{Set}$ is cartesian closed if and only if $*$ is the minimum t-norm. This extends the classification for continuous t-norms established in the first part of this work to the whole left-continuous setting.
\end{abstract}

\maketitle

\section{Introduction}

Building on the theory of frame-valued sets initiated by Higgs \cite{Higgs1970,Higgs1984} and Fourman--Scott \cite{Fourman1979}, H\"ohle and his collaborators \cite{Hoehle1992,Hoehle1995a,Hoehle1995b,Hoehle2011a,Pu2012} developed the theory of quantale-valued sets, which associates with a unital involutive quantale $\sQ$ a category 
\[
\QSet
\]
of sets valued in $\sQ$. From the perspective of enriched category theory \cite{Stubbe2014}, these are symmetric categories enriched in the quantaloid $\sD\sQ$; their morphisms are left adjoint $\sD\sQ$-distributors. When $\sQ$ is a frame, $\QSet$ is a topos \cite{Fourman1979}. For a general quantale, however, $\QSet$ need not be a topos \cite{Hu2025}, and it can even fail to be cartesian closed \cite{Shen2026}.

In this paper, the table $[0,1]_*$ of truth-values is the unit interval equipped with a left-continuous t-norm. The first part of this work \cite{Shen2026} established that, when $*$ is continuous, $\RSet$ is cartesian closed if and only if $*$ is the minimum t-norm.

The left-continuous case cannot be treated as a formal extension of the continuous one. Indeed, Lai and Luo \cite{Lai2026a} show that, after continuity is dropped, the category $[0,1]_*\text{-}\Cat$ of categories enriched in the unital quantale $([0,1],*,1)$ can be cartesian closed for certain non-minimum left-continuous t-norms; see also the general quantaloid characterization of Stubbe and Yu \cite{Stubbe2026}. The category $[0,1]_*\text{-}\Cat$ is a global-enrichment setting, whereas $\RSet$ allows varying type data. This contrast is not explained merely by the availability of variable types: $\RSet$ has symmetric $\sD[0,1]_*$-categories with variable types as objects and left adjoint $\sD[0,1]_*$-distributors as morphisms, whereas $[0,1]_*\text{-}\Cat$ has $[0,1]_*$-categories and $[0,1]_*$-functors. We prove that $\RSet$ is cartesian closed only for the minimum t-norm; in particular, the non-minimum cases in the classification of Lai and Luo do not yield exceptions here.

Our proof differs substantially from the one in Part~I \cite{Shen2026}. Rather than using the structural description of continuous t-norms to select a local configuration and then deriving a full formula for the homs of a hypothetical exponential, we attach a canonical local obstruction to each non-idempotent $c\in[0,1]$. Namely, we consider the idempotent floor
\[
a_c=\bigvee\{e\leq c\mid e*e=e\}.
\]
The preservation of suprema by $*$ makes $a_c$ idempotent, and hence $a_c<c$. From the pair $a_c<c$ we construct two small variable-type test objects. If a suitable exponential existed, their Cauchy completions would force two lower bounds on homs in the exponential. Evaluation gives an incompatible upper bound. Thus the argument applies throughout the left-continuous setting without a full calculation of the homs of a hypothetical exponential.

The paper is organized as follows. Section \ref{real-valued-sets} recalls real-valued sets and their distributor presentation. In Section \ref{cartesian-closedness-of-real-valued-sets}, we first collect the required facts about Cauchy completion, then prove the canonical obstruction and the cartesian-closedness characterization, and finally compare it with the result of Lai and Luo for $[0,1]_*\text{-}\Cat$.

\section{Real-valued sets}
\label{real-valued-sets}

Throughout this paper, we consider $[0,1]_*$ as the table of truth-values, where $*$ is a left-continuous t-norm on the unit interval $[0,1]$. More generally, a binary operation $*$ on an interval $[a,b]$ is a \emph{left-continuous t-norm} \cite{Klement2000,Klement2004b,Alsina2006} if
\begin{itemize}
\item $([a,b],*,b)$ is a commutative monoid,
\item $p*\bigvee\limits_{i\in I}q_{i}=\bigvee\limits_{i\in I}(p*q_{i})$ and $(\bigvee\limits_{i\in I}p_{i})*q=\bigvee\limits_{i\in I}(p_{i}*q)$ for all $p,q,p_i,q_i\in[a,b]$ $(i\in I)$.
\end{itemize}
We write $[a,b]_*$ for the resulting t-norm structure.


Note that for each $p\in[a,b]$, there is a Galois connection
\[(p*-)\dv(p\ra -)\colon[a,b]\to[a,b]\]
satisfying
\[p*q\leq r\iff q\leq p\ra r\]
for all $p,q,r\in[a,b]$. 

Let $p\in[a,b]$. We say that $p$ is \emph{idempotent} if $p*p=p$.

% The following continuous examples were used in Part I.
% \begin{exmp} \label{MLP-def}
% The following continuous t-norms on the unit interval $[0,1]$ are the most prominent ones:
% \begin{enumerate}[label=(\arabic*)]
% \item \label{MLP-def:M} The \emph{minimum t-norm} $[0,1]_{\wedge}$ is given by the meet of real numbers, in which every $q\in(0,1)$ is idempotent, and
% \[p\ra q=\begin{cases}
% 1 & \text{if}\ p\leq q,\\
% q & \text{if}\ p>q.
% \end{cases}\]
% \item \label{MLP-def:P} The \emph{product t-norm} $[0,1]_{\times}$ is given by the usual multiplication of real numbers, in which every $q\in(0,1)$ is non-idempotent, and
% \[p\ra q=1\wedge\dfrac{q}{p}.\]
% \item \label{MLP-def:L} The \emph{{\L}ukasiewicz t-norm} $[0,1]_{*_{\L}}$ is given by $p*_{\text{\L}} q=0\vee(p+q-1)$ for all $p,q\in[0,1]$, in which every $q\in(0,1)$ is non-idempotent, and
% \[p\ra q=1\wedge(1-p+q).\]
% \end{enumerate}
% \end{exmp}

\begin{exmp} \label{MLP-def}
The \emph{nilpotent minimum t-norm} $*_{\mathrm{NM}}$ on $[0,1]$ is a left-continuous but non-continuous t-norm, which is given by
\[
p*_{\mathrm{NM}}q=
\begin{cases}
0 & \text{if}\ p+q\leq1,\\
p\wedge q & \text{if}\ p+q>1.
\end{cases}
\]
Its residual is
\[
p\ra q=
\begin{cases}
1 & \text{if}\ p\leq q,\\
(1-p)\vee q & \text{if}\ p>q.
\end{cases}
\]
\end{exmp}

In fact,
\[([a,b],*,b)\]
is a commutative unital \emph{quantale} \cite{Rosenthal1990}. 

Given $q,u\in[a,b]$, we say that $u$ is \emph{divisible by $q$} if
\[u=q*(q\ra u).\]

\begin{prop} \label{divisible} 
In every left-continuous t-norm $[a,b]_*$,
\begin{equation}\label{divisible:u-v-q}
v*(q\ra u)=(q\ra v)*u
\end{equation}
whenever $u,v\in[a,b]$ are divisible by $q$.
\end{prop}

\begin{proof}
Since $u=q*(q\ra u)$ and $v=q*(q\ra v)$, associativity and commutativity give
\[
v*(q\ra u)=q*(q\ra v)*(q\ra u)=(q\ra v)*u.\qedhere
\]
\end{proof}

Following the definition of \emph{quantale-valued set} \cite{Hoehle1992,Hoehle1995b,Hoehle2011a,Pu2012,Lai2020}, by a \emph{$[0,1]_*$-set} (or a \emph{$[0,1]_*$-valued set}) we mean a set $X$ equipped with a map
\[\al\colon X\times X\to[0,1],\]
such that
\begin{enumerate}[label=(S\arabic*)]
\item \label{S1} $\al(x,y)=\al(x,x)*(\al(x,x)\ra\al(x,y))=\al(y,y)*(\al(y,y)\ra\al(x,y))$,
\item \label{S2} $\al(x,y)=\al(y,x)$,
\item \label{S3} $\al(y,z)*(\al(y,y)\ra\al(x,y))\leq\al(x,z)$
\end{enumerate}
for all $x,y,z\in X$. Here 
\[\al(x,y)\] 
is understood as the truth-value of the statement that \emph{$x$ is equal to $y$}. 

In particular, it necessarily holds that
\begin{equation}\label{al(x,y) leq al(x,x) wedge al(y,y)}
\al(x,y)\leq\al(x,x)\wedge\al(y,y).
\end{equation}
For $x\in X$, we call the value $\al(x,x)$ the \emph{type}, or \emph{extent}, of $x$, which may be understood as the \emph{extent of existence} of $x$.

Let $(X,\al)$ be a $[0,1]_*$-set. For $x,y\in X$, we write $x\cong y$ if
\begin{equation} \label{x-cong-y-def}
\al(x,x)=\al(y,y)=\al(x,y),
\end{equation}
and we say that $(X,\al)$ is \emph{separated} if 
\[x\cong y\iff x=y.\]
Moreover, we denote by
\begin{equation} \label{X_q-def}
X_q=\{x\in X\mid\al(x,x)=q\}
\end{equation}
the slice of $X$ consisting of elements of type $q$.

A \emph{distributor} 
\[\phi\colon (X,\al)\oto (Y,\be)\]
of $[0,1]_*$-sets is a function
\[\phi\colon X\times Y\to[0,1]\]
such that
\begin{enumerate}[label=(M\arabic*)]
\item \label{M1} $\phi(x,y)=\al(x,x)*(\al(x,x)\ra\phi(x,y))=\be(y,y)*(\be(y,y)\ra\phi(x,y))$,
\item \label{M2} $(\be(y,y)\ra\be(y,y'))*\phi(x,y)*(\al(x,x)\ra\al(x',x))\leq\phi(x',y')$
\end{enumerate}
for all $x,x'\in X$, $y,y'\in Y$; it necessarily holds that
\begin{equation}\label{phi(x,y) leq al(x,x) wedge be(y,y)}
\phi(x,y)\leq\al(x,x)\wedge\be(y,y).
\end{equation}
The \emph{opposite} of $\phi$ is the distributor
\[\phi^{\circ}\colon(Y,\be)\oto(X,\al),\quad \phi^{\circ}(y,x)=\phi(x,y).\]
Its composite
\[\psi\circ\phi\colon(X,\al)\oto(Z,\ga)\]
with another distributor $\psi\colon(Y,\be)\oto(Z,\ga)$ is given by
\begin{equation} \label{psi-circ-phi-def}
(\psi\circ\phi)(x,z)=\bv_{y\in Y}\psi(y,z)*(\be(y,y)\ra\phi(x,y)),
\end{equation}
with
\begin{equation} \label{al-id}
\al\colon(X,\al)\oto(X,\al)
\end{equation}
serving as the identity distributor for this composition. The category of $[0,1]_*$-sets and their distributors is denoted by
\[\RDist.\]
$\RDist$ is a \emph{quantaloid} \cite{Rosenthal1996,Stubbe2005}; that is, each hom-set is a complete lattice, the order on distributors is pointwise, and the composition preserves suprema in each variable, i.e.,
\[\psi\circ\Big(\bigvee_{i\in I} \phi_i\Big)=\bigvee_{i\in I}\psi\circ \phi_i\quad\text{and}\quad\Big(\bigvee_{i\in I}\psi_i\Big)\circ \phi=\bigvee_{i\in I}\psi_i\circ \phi\]
for all distributors $\phi,\phi_i\colon (X,\al)\oto(Y,\be)$, $\psi,\psi_i\colon(Y,\be)\oto (Z,\ga)$ $(i\in I)$. The corresponding right adjoints induced by the composition maps
\[(-\circ \phi)\dashv(-\lda \phi)\colon\RDist(X,Z)\to\RDist(Y,Z)\]
and
\[(\psi\circ -)\dashv(\psi\rda -)\colon\RDist(X,Z)\to\RDist(X,Y)\]
satisfy
\[\psi\circ \phi\leq\xi\iff \psi\leq\xi\lda \phi\iff \phi\leq \psi\rda\xi\]
for any $\phi\colon (X,\al)\oto (Y,\be)$, $\psi\colon (Y,\be)\oto (Z,\gamma)$, $\xi\colon (X,\al)\oto (Z,\gamma)$. 

A \emph{morphism} 
\[\phi\colon(X,\al)\oto(Y,\be)\] 
of $[0,1]_*$-sets is a left adjoint distributor, whose right adjoint is necessarily its opposite (see \cite[Proposition 3.5.3]{Heymans2010} and \cite[Theorem 3.7]{Heymans2011}); equivalently, it is a distributor satisfying
\begin{enumerate}[label=(M\arabic*)]
\setcounter{enumi}{2}
\item \label{M3} $\phi\circ\phi^{\circ}\leq\be$,
\item \label{M4} $\al\leq\phi^{\circ}\circ\phi$.
\end{enumerate}
%Each identity distributor is a morphism, and the composite of two morphisms, viewed as distributors, is also a morphism. 
The category of $[0,1]_*$-sets and their morphisms is denoted by
\[\RSet.\]

We may also consider a \emph{monotone function}
\[f\colon(X,\al)\to(Y,\be),\]
which is a map $f\colon X\to Y$ such that
\begin{equation} \label{monotone-function-def}
\al(x,x)=\be(fx,fx)\quad\text{and}\quad\al(x,x')\leq\be(fx,fx')
\end{equation}
for all $x,x'\in X$. The category of $[0,1]_*$-sets and monotone functions is denoted by
\[\RSetM.\]

\begin{rem}
Every left-continuous t-norm $[0,1]_*$ gives rise to a quantaloid $\sD[0,1]_*$ \cite{Hoehle2011,Pu2012,Stubbe2014,Lai2020}. Its objects are the elements of $[0,1]$, and, for $p,q\in[0,1]$,
\[\sD[0,1]_*(p,q)=\{u\in[0,1]\mid u\text{ is divisible by both }p\text{ and }q\}.\]
For $u\in\sD[0,1]_*(p,q)$ and $v\in\sD[0,1]_*(q,r)$, the composite is $v\circ u=v*(q\ra u)$; the identity $\sD[0,1]_*$-arrow of $q$ is $q$ itself, and each hom-set is ordered as a subset of $[0,1]$.

The $[0,1]_*$-sets and distributors defined above are precisely symmetric $\sD[0,1]_*$-categories and $\sD[0,1]_*$-distributors, respectively. Accordingly, a morphism is precisely a left adjoint $\sD[0,1]_*$-distributor, and a monotone function is precisely a $\sD[0,1]_*$-functor. Thus, $\RSet$ is the category of symmetric $\sD[0,1]_*$-categories and left adjoint $\sD[0,1]_*$-distributors, while $\RSetM$ is the category of symmetric $\sD[0,1]_*$-categories and $\sD[0,1]_*$-functors.
\end{rem}

Every monotone function $f\colon(X,\al)\to(Y,\be)$ induces a morphism 
\begin{equation} \label{f-graph-def}
f_{\nat}\colon(X,\al)\oto(Y,\be),\quad f_{\nat}(x,y)=\be(fx,y)
\end{equation}
of $[0,1]_*$-sets, called the \emph{graph} of $f$, whose opposite
\[f^{\nat}\colon(Y,\be)\oto(X,\al),\quad f^{\nat}(y,x)=\be(y,fx),\]
is called the \emph{cograph} of $f$. Obviously, for every $[0,1]_*$-set $(X,\al)$, the identity function $1_X$ is monotone, and
\[\al=(1_X)_{\nat}=1_X^{\nat}.\]
Hence, in order to simplify the notation, we abbreviate a $[0,1]_*$-set $(X,\al)$ to $X$, and write $1_X^{\nat}(x,y)$ instead of $\al(x,y)$ if no confusion arises. 

For each $q\in[0,1]$, we have a one-element $[0,1]_*$-set $\{q\}$ with 
\[1_{\{q\}}^{\nat}(q,q)=q.\]

For each distributor $\phi\colon X\oto Y$,
\[\phi(x,-)\colon\{1_X^{\nat}(x,x)\}\oto Y\]
and
\[\phi(-,y)\colon X\oto \{1_Y^{\nat}(y,y)\}\]
are distributors for all $x\in X$, $y\in Y$. In particular, $\phi(x,y)$ can be considered as a distributor from $\{1_X^{\nat}(x,x)\}$ to $\{1_Y^{\nat}(y,y)\}$, so \ref{M2} can be written as
\begin{enumerate}[label=(M\arabic*$^\prime$)]
\setcounter{enumi}{1}
\item \label{M2'} $1_Y^{\nat}(y,y')\circ\phi(x,y)\circ 1_X^{\nat}(x',x)\leq\phi(x',y')$,
\end{enumerate}
\ref{S3} can be written as
\begin{enumerate}[label=(S\arabic*$^\prime$)]
\setcounter{enumi}{2}
\item \label{S3'} $1_X^{\nat}(y,z)\circ 1_X^{\nat}(x,y)\leq 1_X^{\nat}(x,z)$,
\end{enumerate}
and the composite distributor \eqref{psi-circ-phi-def} can be written as
\begin{equation} \label{psi-circ-phi-def'}
(\psi\circ\phi)(x,z)=\bv_{y\in Y}\psi(y,z)\circ\phi(x,y).
\end{equation}


\begin{lem} \label{phi-p-q}
For $p,q\in[0,1]$, the following statements are equivalent:
\begin{enumerate}[label={\rm(\roman*)}]
\item \label{phi-p-q:phi} There exists a morphism $\phi\colon\{p\}\oto\{q\}$ of one-element $[0,1]_*$-sets.
\item \label{phi-p-q:p} $p$ is divisible by $q$ and satisfies $p=p*(q\ra p)$.
\end{enumerate}
In this case, it necessarily holds that $\phi(p,q)=p$. Therefore, a morphism between one-element $[0,1]_*$-sets, when it exists, will be denoted by
\[\overline{p}\colon\{p\}\oto\{q\}.\]
When the codomain needs to be indicated, we write $\overline{p}_q$ instead of $\overline{p}$.
\end{lem}
\begin{proof}
\ref{phi-p-q:phi}$\implies$\ref{phi-p-q:p}: Suppose that $\phi(p,q)=r$. Then, by \ref{M1} and \ref{M4} we have
\begin{equation} \label{phi-p-q:M134}
r\leq p\wedge q\quad\text{and}\quad p\leq r*(q\ra r).
\end{equation}
Thus
\[p\leq r*(q\ra r)\leq r\leq p,\] 
and consequently $r=p$. By \ref{M1},
\[p=r=q*(q\ra r)=q*(q\ra p),\]
so $p$ is divisible by $q$. Moreover,
\[p\leq r*(q\ra r)=p*(q\ra p)\leq p;\]
hence
\begin{equation} \label{p=p*(q-ra-p)}
p=p*(q\ra p).
\end{equation}
\ref{phi-p-q:p}$\implies$\ref{phi-p-q:phi}: In this case, it is easy to see that $\phi(p,q)=p$ satisfies \ref{M1}--\ref{M4}:
\begin{itemize}
\item \ref{M1} $p=p*(p\ra p)=q*(q\ra p)$.
\item \ref{M2} $(q\ra q)*p*(p\ra p)\leq p$.
\item \ref{M3} $p*(p\ra p)\leq q$.
\item \ref{M4} $p\leq p*(q\ra p)$.
\end{itemize}
Therefore, $\phi\colon\{p\}\oto\{q\}$ is a morphism of one-element $[0,1]_*$-sets.
\end{proof} 

\section{Cartesian closedness of the category of real-valued sets}
\label{cartesian-closedness-of-real-valued-sets}

A \emph{singleton} on a $[0,1]_*$-set $X$ is a morphism $\lam\colon\{p\}\oto X$ whose domain is a one-element $[0,1]_*$-set. A $[0,1]_*$-set is \emph{Cauchy complete} precisely when every singleton is represented by an element.

\begin{prop} \label{Cauchy-complete-def} (See \cite[Proposition 7.1]{Stubbe2005}.)
A $[0,1]_*$-set $X$ is Cauchy complete if and only if for each singleton $\lam\colon\{p\}\oto X$, there exists $x\in X$ such that
\[\lam=1_X^{\nat}(x,-).\]
In this case, necessarily $1_X^{\nat}(x,x)=p$.
\end{prop}

For every $[0,1]_*$-set $X$, let
\[\CdX\coloneqq\bigcup_{q\in[0,1]}\RSet(\{q\},X)\]
be the set of all singletons on $X$. Its hom is
\begin{equation} \label{CdX-hom}
1_{\CdX}^{\nat}(\lam,\lam')={\lam'}^{\circ}\circ\lam.
\end{equation}
The $[0,1]_*$-set $\CdX$ is separated and Cauchy complete \cite[Proposition 7.12]{Stubbe2005}; it is called the \emph{Cauchy completion} of $X$. For a singleton $\lam\colon\{q\}\oto X$, one has
\begin{equation} \label{CdX-type}
1_{\CdX}^{\nat}(\lam,\lam)=q.
\end{equation}
The canonical monotone function
\begin{equation} \label{FYd-def}
\Fyd_X\colon X\to\CdX,\quad \Fyd_X x=1_X^{\nat}(x,-)
\end{equation}
is the unit of the Cauchy-completion reflection.

The following two elementary descriptions are obtained by the same formal calculations as in \cite[Examples 3.3 and 3.4]{Shen2026}. The difference is that Lemma \ref{phi-p-q} replaces the continuous-case description of morphisms between one-element $[0,1]_*$-sets:

\begin{exmp} \label{Cdq-exmp}
By Lemma \ref{phi-p-q}, the Cauchy completion of the one-element $[0,1]_*$-set $\{q\}$ is
\begin{equation} \label{Cdq}
\sCd\{q\}=\{\overline{p}\mid p\text{ is divisible by }q\text{ and }p=p*(q\ra p)\}.
\end{equation}
For $\overline{p},\overline{p'}\in\sCd\{q\}$,
\begin{equation} \label{Cdp-hom}
1_{\sCd\{q\}}^{\nat}(\overline{p},\overline{p'})=p'*(q\ra p).
\end{equation}
\end{exmp}

\begin{exmp} \label{Cdq-x-y-exmp}
Let $X=\{x,y\}$ be a two-element $[0,1]_*$-set. Every singleton on $X$ has a presentation
\[
1_X^{\nat}(z,-)\circ\overline{p}\colon\{p\}\oto X,
\]
where $z\in\{x,y\}$ and $\overline{p}\colon\{p\}\oto\{1_X^{\nat}(z,z)\}$ is a morphism. Hence
\[
\sCd X=\{1_X^{\nat}(z,-)\circ\overline{p}\mid z\in\{x,y\},\ \overline{p}\colon\{p\}\oto\{1_X^{\nat}(z,z)\}\}.
\]
\end{exmp}

Let $\RCcSetM$ be the full subcategory of $\RSetM$ consisting of separated Cauchy complete $[0,1]_*$-sets. It is reflective in $\RSetM$, with reflector $\sCd$ (see \cite[Proposition 7.14]{Stubbe2005}). The following standard equivalence can be proved by the same argument as \cite[Proposition 5.7(2)]{Pu2012}; the argument remains valid for left-continuous t-norms in the present setting (cf. \cite{Stubbe2005,Hoehle2011a}). It allows us to study the cartesian closedness of $\RSet$ in $\RCcSetM$.

\begin{prop} \label{RSet-RCcSetM}
The category $\RSet$ is equivalent to $\RCcSetM$.
\end{prop}


Recall that an object $Y$ of a category with finite products is \emph{exponentiable} if the functor $-\times Y$ has a right adjoint $(-)^Y$. Thus, if $Z^Y$ exists, there is an evaluation
\[\ev_Z\colon Z^Y\times Y\to Z\]
such that every monotone function $K\colon A\times Y\to Z$ has a unique transpose
\[\sE K\colon A\to Z^Y\]
with $\ev_Z\circ(\sE K\times1_Y)=K$.

The category $\RCcSetM$ has finite products. For $A,B\in\RCcSetM$,
\begin{equation} \label{X-times-Y-def}
A\times B=\{(x,y)\mid x\in A,\ y\in B,\ 1_A^{\nat}(x,x)=1_B^{\nat}(y,y)\}
\end{equation}
and
\[
1_{A\times B}^{\nat}\left((x,y),(x',y')\right)=1_A^{\nat}(x,x')\wedge1_B^{\nat}(y,y').
\]

We first isolate the only consequence of the exponential adjunction that will be used below.

\begin{lem} \label{Z^Y-hom}
Let $B$ be a $[0,1]_*$-set, let $Y,Z\in\RCcSetM$, and suppose that $Z^Y$ exists in $\RCcSetM$. For every monotone function
\[
K\colon\sCd B\times Y\to Z,
\]
put $k=\sE K\colon\sCd B\to Z^Y$. Then
\[
1_B^{\nat}(b,b')\leq1_{Z^Y}^{\nat}\bigl(k(\Fyd_Bb),k(\Fyd_Bb')\bigr)
\]
for all $b,b'\in B$.
\end{lem}

\begin{proof}
The composite $k\circ\Fyd_B\colon B\to Z^Y$ is monotone. The assertion follows immediately from \eqref{monotone-function-def}.
\end{proof}

\begin{lem} \label{sup-idem}
For $c\in[0,1]$, put
\[
a_c=\bigvee\{e\leq c\mid e*e=e\}.
\]
Then $a_c$ is idempotent. Consequently, if $c$ is non-idempotent, then $a_c<c$, and every idempotent element below $c$ is at most $a_c$.
\end{lem}

\begin{proof}
Let $E_c=\{e\leq c\mid e*e=e\}$. Since $*$ preserves suprema in each variable,
\[
a_c*a_c=\bigvee_{e,e'\in E_c}e*e'.
\]
For $e,e'\in E_c$, one has $e*e'\leq e\leq a_c$, whereas the diagonal terms give $e=e*e\leq a_c*a_c$. Hence $a_c*a_c=a_c$. If $a_c=c$, then $c$ is idempotent, so a non-idempotent $c$ satisfies $a_c<c$. The final assertion is immediate from the definition of $a_c$.
\end{proof}

\begin{prop} \label{RCcSetM-cc}
The category $\RCcSetM$ is cartesian closed if and only if $*$ is the minimum t-norm on $[0,1]$.
\end{prop}

\begin{proof}
If $*$ is the minimum t-norm, then $\RSet$ is the topos of $[0,1]$-valued sets over the frame $[0,1]$ \cite{Hoehle1992}. It is therefore cartesian closed, and so is the equivalent category $\RCcSetM$ by Proposition \ref{RSet-RCcSetM}.

Conversely, suppose that $*$ is not the minimum t-norm. There is then a non-idempotent $c\in[0,1]$: indeed, if every element were idempotent, then $p*q=p\wedge q$ for all $p,q\in[0,1]$. By Lemma \ref{sup-idem}, $a_c<c$.

Let $X=\{x,x'\}$ be the $[0,1]_*$-set with
\[
1_X^{\nat}(x,x)=1_X^{\nat}(x',x')=1,\qquad
1_X^{\nat}(x,x')=a_c.
\]
Put $Y=\sCd\{1\}$ and $Z=\sCd X$. An element of $Y$ is of the form $\overline{p}_1\colon\{p\}\oto\{1\}$, where $p$ is idempotent. If $(\overline{p}_c,\overline{p}_1)$ belongs to $\sCd\{c\}\times Y$, then $p$ is an idempotent element below $c$, and hence $p\leq a_c$. Since $a_c$ and $p$ are idempotent, this also gives $a_c*p=p$.

For idempotent $p,q$, direct use of \eqref{CdX-hom} gives
\[
1_Z^{\nat}\left(1_X^{\nat}(x,-)\circ\overline{p}_1,1_X^{\nat}(x,-)\circ\overline{q}_1\right)=p*q
\]
and
\[
1_Z^{\nat}\left(1_X^{\nat}(x,-)\circ\overline{p}_1,1_X^{\nat}(x',-)\circ\overline{q}_1\right)=a_c*p*q.
\]
Thus the following maps are monotone:
\[
F\colon\sCd\{1\}\times Y\to Z,\qquad
F(\overline{p}_1,\overline{p}_1)=1_X^{\nat}(x,-)\circ\overline{p}_1,
\]
\[
G\colon\sCd\{c\}\times Y\to Z,\qquad
G(\overline{p}_c,\overline{p}_1)=1_X^{\nat}(x,-)\circ\overline{p}_1,
\]
and
\[
H\colon\sCd\{1\}\times Y\to Z,\qquad
H(\overline{p}_1,\overline{p}_1)=1_X^{\nat}(x',-)\circ\overline{p}_1.
\]
For $G$, the preceding observation and the displayed hom calculations show monotonicity by the product formula.

Suppose that $Z^Y$ exists. Since $Z^Y$ is Cauchy complete, the exponential adjunction and the Cauchy-completion reflection give natural bijections (cf. \cite[(4.vi)--(4.viii)]{Shen2026})
\[
Z^Y_q\cong\RCcSetM(\sCd\{q\}\times Y,Z)
\]
for $q\in[0,1]$. Let $f\in Z^Y_1$, $g\in Z^Y_c$, and $h\in Z^Y_1$ be the elements corresponding to $F$, $G$, and $H$, respectively.

Let $B=\{b_0,b_1\}$ be the $[0,1]_*$-set with
\[
1_B^{\nat}(b_0,b_0)=1,\qquad
1_B^{\nat}(b_1,b_1)=c,\qquad
1_B^{\nat}(b_0,b_1)=c.
\]
Define a map $K\colon\sCd B\times Y\to Z$ by
\[
K\left(1_B^{\nat}(b_i,-)\circ\lam,\overline{p}_1\right)
=1_X^{\nat}(x,-)\circ\overline{p}_1,
\]
where $\lam\colon\{p\}\oto\{1_B^{\nat}(b_i,b_i)\}$ is a morphism. By Example \ref{Cdq-x-y-exmp}, this defines $K$ on all of $\sCd B\times Y$. If a point has two such presentations, the resulting value of $K$ is the same for both presentations. Moreover, the hom in the source product is at most $p*q$, which is exactly the hom between the corresponding images in $Z$. Hence $K$ is monotone.

For $i=0,1$, the singleton $1_B^{\nat}(b_i,-)$ induces a monotone function
\[
j_i\colon\sCd\{1_B^{\nat}(b_i,b_i)\}\to\sCd B,\qquad
j_i(\lam)=1_B^{\nat}(b_i,-)\circ\lam.
\]
The restrictions of $K$ along $j_0\times1_Y$ and $j_1\times1_Y$ are $F$ and $G$, respectively. Under the natural bijections above, naturality with respect to $j_0$ and $j_1$ identifies the corresponding values of the transpose $\sE K$ with $f$ and $g$; explicitly,
\[
(\sE K)(\Fyd_B b_0)=f,\qquad (\sE K)(\Fyd_B b_1)=g.
\]
Lemma \ref{Z^Y-hom} therefore gives
\[
c\leq1_{Z^Y}^{\nat}(f,g).
\]
The reverse inequality follows from the types of $f$ and $g$, so
\[
1_{Z^Y}^{\nat}(f,g)=c.
\]

Similarly, let $B'=\{b'_0,b'_1\}$ be the $[0,1]_*$-set with
\[
1_{B'}^{\nat}(b'_0,b'_0)=c,\qquad
1_{B'}^{\nat}(b'_1,b'_1)=1,\qquad
1_{B'}^{\nat}(b'_0,b'_1)=c.
\]
Define $K'\colon\sCd B'\times Y\to Z$ by
\[
K'\left(1_{B'}^{\nat}(b'_0,-)\circ\lam,\overline{p}_1\right)
=1_X^{\nat}(x,-)\circ\overline{p}_1
\]
and
\[
K'\left(1_{B'}^{\nat}(b'_1,-)\circ\lam,\overline{p}_1\right)
=1_X^{\nat}(x',-)\circ\overline{p}_1.
\]
If a point admits both presentations used in the definition of $K'$, then its type is an idempotent below $c$, hence is at most $a_c$; the two prescriptions for $K'$ therefore give the same element of $Z$, because $a_c*p=p$. In a mixed case, the parameter of a singleton presented through $b'_0$ is likewise at most $a_c$; consequently the mixed hom in $Z$ is $a_c*p*q=p*q$. Since the hom in the source product is at most $p*q$, the map $K'$ is monotone.

For $i=0,1$, define
\[
j'_i\colon\sCd\{1_{B'}^{\nat}(b'_i,b'_i)\}\to\sCd B',\qquad
j'_i(\lam)=1_{B'}^{\nat}(b'_i,-)\circ\lam.
\]
The restrictions of $K'$ along $j'_0\times1_Y$ and $j'_1\times1_Y$ are $G$ and $H$, respectively. The same naturality argument gives
\[
(\sE K')(\Fyd_{B'} b'_0)=g,\qquad (\sE K')(\Fyd_{B'} b'_1)=h.
\]
Hence Lemma \ref{Z^Y-hom} gives
\[
c\leq1_{Z^Y}^{\nat}(g,h).
\]
The reverse inequality follows from the types of $g$ and $h$, so
\[
1_{Z^Y}^{\nat}(g,h)=c.
\]

Let $\overline{1}_1$ denote the type-$1$ element of $Y$. The monotonicity of evaluation gives
\[
1_{Z^Y}^{\nat}(f,h)\leq1_Z^{\nat}\left(\ev_Z(f,\overline{1}_1),\ev_Z(h,\overline{1}_1)\right).
\]
The defining property of the transposes and the displayed calculation of the hom in $Z$ identify the right-hand side with
\[
1_Z^{\nat}\left(1_X^{\nat}(x,-),1_X^{\nat}(x',-)\right)=a_c.
\]
On the other hand, $1_{Z^Y}^{\nat}(g,g)=c$, and axiom \ref{S3} applied to $f,g,h$ yields
\[
c*(c\ra c)\leq1_{Z^Y}^{\nat}(f,h).
\]
This is impossible because $c*(c\ra c)=c>a_c$. Hence $Z^Y$ does not exist, and $\RCcSetM$ is not cartesian closed.
\end{proof}

The main result follows at once from Propositions \ref{RSet-RCcSetM} and \ref{RCcSetM-cc}:

\begin{thm} \label{QSet-topos-frame}
The category $\RSet$ is cartesian closed if and only if $*$ is the minimum t-norm on $[0,1]$.
\end{thm}

\begin{rem}

Theorem \ref{QSet-topos-frame} contrasts with the result of Lai and Luo \cite{Lai2026a}; see also the general quantaloid characterization of Stubbe and Yu \cite{Stubbe2026}. In $[0,1]_*\text{-}\Cat$, enrichment is global: all hom-values lie in the single quantale $[0,1]_*$, so there is no varying type data. Lai and Luo characterize the left-continuous t-norms for which $[0,1]_*\text{-}\Cat$ is cartesian closed; their classification includes t-norms different from the minimum t-norm.

By contrast, $\RSet$ has symmetric $\sD[0,1]_*$-categories with variable types as objects and left adjoint $\sD[0,1]_*$-distributors as morphisms. The two test objects $B$ and $B'$ in the proof of Proposition \ref{RCcSetM-cc} use the non-idempotent type $c$ alongside type $1$, a configuration unavailable in the global-enrichment setting of $[0,1]_*\text{-}\Cat$. Thus none of the non-minimum cases in the classification of Lai and Luo yields a cartesian closed $\RSet$. This contrast is not attributable to variable types alone: symmetry and the choice of left adjoint distributors are also part of the $\RSet$ framework.
\end{rem}

\begin{cor}
The category $\RSet$ is a topos if and only if $*$ is the minimum t-norm on $[0,1]$.
\end{cor}

\begin{proof}
If $*$ is the minimum t-norm, then $\RSet$ is a topos \cite{Hoehle1992}. Conversely, every topos is cartesian closed, so the conclusion follows from Theorem \ref{QSet-topos-frame}.
\end{proof}

% \subsection*{Acknowledgment}

% The authors acknowledge the support of National Natural Science Foundation of China (No. 12071319).

\section*{Acknowledgements}

During the preparation of this manuscript, the authors used ChatGPT to assist with manuscript organization, language editing, mathematical consistency checking, and submission-readiness review. All AI-assisted suggestions were critically reviewed and, where appropriate, edited and verified by the authors, who remain fully responsible for the content of the manuscript.

\section*{Declarations}

\subsection*{Ethical approval}
Not applicable.

\subsection*{Competing interests}
The authors declare that they have no competing interests.

\subsection*{Authors' contributions}
All authors contributed equally.

\subsection*{Availability of data and materials}
Not applicable.

\subsection*{Funding}

The authors acknowledge support from the National Natural Science Foundation of China (Grant No. 12671561).

\begin{thebibliography}{10}

\bibitem{Alsina2006}
Alsina, C., Frank, M.J., Schweizer, B.: Associative Functions: Triangular Norms
  and Copulas.
\newblock World Scientific, Singapore (2006).

\bibitem{Fourman1979}
Fourman, M.P., Scott, D.S.: Sheaves and logic.
\newblock In: M.P. Fourman, C.J. Mulvey, D.S. Scott (eds.) Applications of
  Sheaves: Proceedings of the Research Symposium on Applications of Sheaf
  Theory to Logic, Algebra, and Analysis, Durham, July 9--21, 1977,
  \emph{Lecture Notes in Mathematics}, vol. 753, pp. 302--401. Springer,
  Berlin--Heidelberg (1979).

\bibitem{Heymans2010}
Heymans, H.: Sheaves on quantales as generalized metric spaces.
\newblock Ph.D. thesis, Universiteit Antwerpen, Belgium (2010).

\bibitem{Heymans2011}
Heymans, H., Stubbe, I.: Symmetry and {Cauchy} completion of
  quantaloid-enriched categories.
\newblock Theory and Applications of Categories \textbf{25}(11), 276--294
  (2011).

\bibitem{Higgs1970}
Higgs, D.: Boolean-valued equivalence relations and complete extensions of
  complete boolean algebras.
\newblock Bulletin of the Australian Mathematical Society \textbf{3}(1), 65--72
  (1970).

\bibitem{Higgs1984}
Higgs, D.: Injectivity in the topos of complete {Heyting} algebra valued sets.
\newblock Canadian Journal of Mathematics \textbf{36}(3), 550--568 (1984).

\bibitem{Hoehle1992}
H{\"o}hle, U.: {$M$}-valued sets and sheaves over integral commutative
  cl-monoids.
\newblock In: S.E. Rodabaugh, E.P. Klement, U.~H\"{o}hle (eds.) Applications of
  Category Theory to Fuzzy Subsets, \emph{Theory and Decision Library},
  vol.~14, pp. 33--72. Springer, Berlin--Heidelberg (1992).

\bibitem{Hoehle1995a}
H{\"o}hle, U.: Commutative, residuated l-monoids.
\newblock In: U.~H\"{o}hle, E.P. Klement (eds.) Non-classical logics and their
  applications to fuzzy subsets: a handbook of the mathematical foundations of
  fuzzy set theory, \emph{Theory and Decision Library}, vol.~32, pp. 53--105.
  Springer, Berlin--Heidelberg (1995).

\bibitem{Hoehle1995b}
H{\"o}hle, U.: Presheaves over {GL}-monoids.
\newblock In: U.~H\"{o}hle, E.P. Klement (eds.) Non-Classical Logics and their
  Applications to Fuzzy Subsets, \emph{Theory and Decision Library}, vol.~32,
  pp. 127--157. Springer, Berlin--Heidelberg (1995).

\bibitem{Hoehle2011a}
H{\"o}hle, U., Kubiak, T.: A non-commutative and non-idempotent theory of
  quantale sets.
\newblock Fuzzy Sets and Systems \textbf{166}, 1--43 (2011).

\bibitem{Hoehle2011}
H{\"o}hle, U., Kubiak, T.: On regularity of sup-preserving maps: generalizing
  {Zarecki{\u{\i}}}'s theorem.
\newblock Semigroup Forum \textbf{83}(2), 313--319 (2011).

\bibitem{Hu2025}
Hu, X., Shen, L.: $\mathsf{Q}\text{-}\mathbf{Set}$ is not generally a topos.
\newblock Fuzzy Sets and Systems \textbf{517}, 109484 (2025).

\bibitem{Klement2000}
Klement, E.P., Mesiar, R., Pap, E.: Triangular Norms, \emph{Trends in Logic},
  vol.~8.
\newblock Springer, Dordrecht (2000).

\bibitem{Klement2004b}
Klement, E.P., Mesiar, R., Pap, E.: Triangular norms. {Position} paper {III}:
  continuous t-norms.
\newblock Fuzzy Sets and Systems \textbf{145}(3), 439--454 (2004).

\bibitem{Lai2026a}
Lai, H., Luo, Q.: On the {Cartesian} closedness of
  {$[0,1]\text{-}\mathsf{Cat}$} and some of its subcategories.
\newblock Theory and Applications of Categories \textbf{45}(8), 289--306
  (2026).

\bibitem{Lai2020}
Lai, H., Shen, L., Tao, Y., Zhang, D.: Quantale-valued dissimilarity.
\newblock Fuzzy Sets and Systems \textbf{390}, 48--73 (2020).

\bibitem{Pu2012}
Pu, Q., Zhang, D.: Preordered sets valued in a {GL}-monoid.
\newblock Fuzzy Sets and Systems \textbf{187}(1), 1--32 (2012).

\bibitem{Rosenthal1990}
Rosenthal, K.I.: Quantales and their Applications, \emph{Pitman research notes
  in mathematics series}, vol. 234.
\newblock Longman, Harlow (1990).

\bibitem{Rosenthal1996}
Rosenthal, K.I.: The Theory of Quantaloids, \emph{Pitman Research Notes in
  Mathematics Series}, vol. 348.
\newblock Longman, Harlow (1996).

\bibitem{Shen2026}
Shen, L., Zhang, J.: Cartesian closedness of the category of real-valued sets,
  {I}.
\newblock Fuzzy Sets and Systems \textbf{531}, 109767 (2026).

\bibitem{Stubbe2005}
Stubbe, I.: Categorical structures enriched in a quantaloid: categories,
  distributors and functors.
\newblock Theory and Applications of Categories \textbf{14}(1), 1--45 (2005).

\bibitem{Stubbe2014}
Stubbe, I.: An introduction to quantaloid-enriched categories.
\newblock Fuzzy Sets and Systems \textbf{256}, 95--116 (2014).

\bibitem{Stubbe2026}
Stubbe, I., Yu, J.: When is $\textsf{Cat}(\mathcal Q)$ Cartesian closed?
\newblock Applied Categorical Structures \textbf{34}(2), 9 (2026).

\end{thebibliography}

\end{document}
